% Chapter 2. Background and related work. No measured values.
\chapter{Background and Related Work}
\label{ch:background}

\section{Cascaded voice agents}
\label{sec:cascade}

A cascaded voice agent turns speech into a reply in stages:
\begin{enumerate}
\item a voice-activity detector decides when the user has started and stopped
  speaking;
\item a speech recognizer transcribes what was said;
\item a language model writes the reply;
\item a synthesizer speaks it.
\end{enumerate}
Each stage waits for the previous one at the turn boundary, so the delay the
user hears is a sum. The endpoint is itself a design choice: a detector waits
for a stretch of silence before declaring the turn over, trading a later
endpoint against cutting the user off. Turn-taking can instead be predicted,
for example by projecting voice activity forward \cite{vap2022}. Human gaps
between turns are short \cite{stivers2009universals}, which is why latency to
the first reply audio, rather than total generation time, is the measure
that matters in conversation.

Cascades remain common on devices because each stage can be chosen, replaced
and run within a small memory budget. Low-latency agent designs, both cascaded
and speech-to-speech, are described in
\cite{ilava2025,chipchat2025,kame2026,streamrag2025,moshirag2026}, and
benchmarks evaluate voice agents and their turn-taking
\cite{fullduplexbench2025,fullduplexbenchv3_2026,voicebench2026,voiceagentbench2025,evabench2026,provoicebench2026,sid2026}.

\section{Speech recognition while the user speaks}
\label{sec:streaming-asr}

Whisper \cite{radford2023whisper} is a recognizer trained with large-scale
weak supervision. Turning it into a real-time transcription system
\cite{machacek2023whisperstreaming} rests on local agreement: committing the
prefix that consecutive hypotheses agree on. Simul-Whisper streams Whisper with
attention-guided chunk decoding and detects words truncated at chunk
boundaries \cite{wang2024simulwhisper}. Other work decides when a streaming
recognizer should emit \cite{x2streaming2026}.

These methods aim at a live transcript. This thesis uses local agreement for a
different end. What it commits before the endpoint is audio that the final
decode no longer has to process, and that shortens the silence after the user
stops.

\section{Thinking while listening}
\label{sec:twl}

\textbf{Speculating the reply in a cascade.} Several systems start the
language model on partial input, through:
\begin{itemize}
\item input-time speculation \cite{predgen2025};
\item semantic triggering with incremental reasoning \cite{ltsvoiceagent2026};
\item endpoint anticipation \cite{epa2026};
\item simultaneous inference \cite{livemind2024};
\item a dual-track streaming response \cite{ddtsr2026}.
\end{itemize}
They differ in when they start, what they draft and how they reconcile the
draft with the final input. All of them spend compute before the endpoint in
the hope that it shortens the time after it.

\textbf{Spoken language models.} Models that take speech in directly
interleave hearing, thinking and speaking
\cite{shanks2025,stitch2026,miniomnireasoner2025,chronothink2026,shih2025thinklisten}.
Full-duplex models remove the turn boundary altogether
\cite{moshi2024,lws2026}. Reasoning models spend more compute at test time
\cite{s1_2025}, which raises the stakes of when that compute can start.

\textbf{The premise.} Whatever the design, its benefit depends on the compute
spent while the user speaks being otherwise idle. On a server
running a reasoner that waits for input, it largely is. On a device where the
recognizer and the reasoner share the hardware, it is not. That is the
setting of this thesis.

\section{Sharing a unified-memory device}
\label{sec:shared}

On an edge system-on-chip, the CPU cores, the GPU and the accelerators share
one physical memory. Workloads that run concurrently interfere through memory
bandwidth, caches, the memory controller's clock and the accelerators
themselves. Interference-aware scheduling and resource control for
concurrent neural networks address this on heterogeneous SoCs and on GPUs
\cite{haxconn2024,aliyun2018bwlock,kim2024interference,kang2026emcgovernor,sgdrc2025,quilt2026},
and in the datacenter \cite{prophet2017,igniter2022}.

\textbf{Transparent huge pages.} Linux can back a process's anonymous memory
with huge pages, which cover more memory per translation entry and so reduce
translation misses. Huge pages need physically contiguous free memory. When
none is available at fault time, the fault falls back to base pages, and the
kernel counts the fallback. Whether the kernel compacts memory to make a huge
page available on demand depends on the system's defragmentation policy. A
process that holds much memory without using it can therefore still affect a
neighbour: it occupies and fragments the pool from which the neighbour's huge
pages come. Chapter~\ref{ch:cost} measures this effect on the recognizer.

\section{Positioning}
\label{sec:positioning}

\begin{table}[t]
\centering\small
\caption{What is known, and what this thesis adds.}
\label{tab:positioning}
\begin{tabular}{p{0.42\textwidth}p{0.48\textwidth}}
\toprule
Known & Added here \\
\midrule
Reply speculation can cut latency when compute is otherwise idle. &
Its cost and value measured on a device where it is not, under one overlap
penalty. \\
Local agreement makes Whisper stream. &
Committing while listening as a way to shorten the final decode, against
reply speculation, on shared hardware. \\
Co-running models interfere on shared SoCs. &
An occupancy tax through huge-page fallback, from a process that does not
run. \\
Speculation needs a trigger. &
A per-turn feasibility rule, its decision log, and a planner showing how
close the rule comes to optimal. \\
\bottomrule
\end{tabular}
\end{table}

Table~\ref{tab:positioning} summarises the position. The mechanisms used here
are not new. What is new is to measure them against each other on hardware
they share, and to give a rule that decides between them from measured
costs.
